\documentclass[border=2pt,tikz]{standalone}
% Shared TikZ/pgfplots style for every figure in the manuscript.
\usepackage[T1]{fontenc}
\usepackage{mathpazo}
\usepackage{amsmath,amssymb}
\usepackage{tikz}
\usepackage{pgfplots}
\usepackage{pgfplotstable}
\pgfplotsset{compat=1.18}
\usetikzlibrary{arrows.meta,positioning,calc,backgrounds,fit,shapes.geometric}
\usepgfplotslibrary{fillbetween}

% --- palette: colour-blind safe, legible in greyscale -------------------
\definecolor{cHyd}{HTML}{1B6CA8}   % hydride / primary blue
\definecolor{cPro}{HTML}{C1492E}   % proton  / primary red
\definecolor{cAcc}{HTML}{7B5AA6}   % accent  / purple
\definecolor{cGold}{HTML}{B8860B}  % accent  / gold
\definecolor{cGrey}{HTML}{6E6E6E}
\definecolor{cInk}{HTML}{1A1A1A}
\definecolor{cBandR}{HTML}{C1492E}
\definecolor{cBandG}{HTML}{4F7942}

% --- base axis ----------------------------------------------------------
\pgfplotsset{
  icig/.style={
    width=68mm, height=52mm,
    scale only axis,
    axis line style={cInk, line width=0.45pt},
    tick style={cInk, line width=0.4pt},
    tick align=outside, tick pos=left,
    label style={font=\small, color=cInk},
    tick label style={font=\footnotesize, color=cInk},
    title style={font=\small\bfseries, color=cInk, align=left},
    legend cell align=left,
    legend style={draw=none, fill=none, font=\scriptsize,
                  inner sep=1pt, row sep=0.5pt},
    grid=both,
    major grid style={cGrey!22, line width=0.28pt},
    minor grid style={cGrey!12, line width=0.2pt},
  },
  % marker shorthands
  ptA/.style={mark=*, mark size=1.7pt, thick,
              mark options={fill=cHyd, draw=white, line width=0.4pt}},
  ptB/.style={mark=*, mark size=1.7pt, thick,
              mark options={fill=cPro, draw=white, line width=0.4pt}},
  hollowB/.style={mark=o, mark size=1.7pt, thick,
              mark options={draw=cPro, line width=0.7pt}},
}

% panel tag, e.g. \panel{a}
\newcommand{\panel}[1]{\textbf{(\textsf{#1})}}

% --- shared notation, matching the manuscript preamble ----------------
\newcommand{\KHT}{K_{\mathrm{HT}}}
\newcommand{\KDT}{K_{\mathrm{DT}}}
\newcommand{\gSC}{\gamma_{\mathrm{SC}}}
\newcommand{\Fz}{F_{0}}

\usepgfplotslibrary{fillbetween}
\begin{document}
\begin{tikzpicture}
% Two-column figure (figure*). Natural width ~150 mm.
\pgfplotsset{
  geo/.style={icig, scale only axis,
    label style={font=\fontsize{8.2}{10}\selectfont, color=cInk},
    tick label style={font=\fontsize{7.4}{9}\selectfont, color=cInk},
    title style={font=\fontsize{8.8}{10.6}\selectfont\bfseries, color=cInk,
                 align=left}},
}
\tikzset{
  ann/.style={font=\fontsize{7.2}{8.6}\selectfont, color=cInk, align=left},
  annc/.style={font=\fontsize{7.2}{8.6}\selectfont, color=cInk, align=center},
  tag/.style={font=\fontsize{7.0}{8.4}\selectfont, align=center,
              inner sep=1.8pt, rounded corners=1pt},
}

% =====================================================================
% (A) the mass-scaling plane
% =====================================================================
\begin{axis}[geo, name=A, width=62mm, height=62mm,
  xmin=0.34, xmax=0.72, ymin=0.95, ymax=3.45,
  xlabel={$X_{\mathrm D}=\ln \KDT$}, ylabel={$X_{\mathrm H}=\ln \KHT$},
  xtick={0.4,0.5,0.6,0.7}, ytick={1,1.5,2,2.5,3}, clip=true,
  title={(\textit{\textbf{A}})\ \ the offset is a vertical distance,\\
         \hspace*{3.1mm}and it can only increase}]

  % region above the semiclassical locus
  \addplot[name path=LOC, cInk, line width=0.9pt] table[x=xd,y=xh]
    {data/fg_locus.dat};
  \path[name path=TOP] (axis cs:0.34,3.45) -- (axis cs:0.72,3.45);
  \addplot[cBandG!12] fill between[of=LOC and TOP];

  % commitment trajectories
  \addplot[cPro, line width=1.1pt] table[x=xd,y=xh] {data/fg_yadh.dat};
  \addplot[cHyd, line width=1.1pt, densely dashed] table[x=xd,y=xh]
    {data/fg_light.dat};
  \addplot[only marks, mark=*, mark size=2.1pt, cPro]
    table[x=xd,y=xh] {data/fg_yadh_obs.dat};
  \addplot[only marks, mark=*, mark size=2.1pt, cHyd]
    table[x=xd,y=xh] {data/fg_light_obs.dat};
  \addplot[only marks, mark=o, mark size=2.4pt, cHyd, line width=0.7pt]
    table[x=xd,y=xh] {data/fg_cross.dat};

  % the offset as an explicit vertical gap at the yeast observation
  \draw[cPro, line width=0.9pt, {Straight Barb[length=1.6pt,width=2.4pt]}-%
        {Straight Barb[length=1.6pt,width=2.4pt]}]
    (axis cs:0.548121,1.835529) -- (axis cs:0.548121,1.964440);

  % annotations, placed clear of the curves and of each other
  \node[ann, anchor=east, color=cPro] at (axis cs:0.536,2.06)
    {$F_{\mathrm{obs}}>0$};
  \draw[cPro!55, line width=0.4pt] (axis cs:0.538,2.03) -- (axis cs:0.547,1.95);
  \node[ann, anchor=west, color=cHyd] at (axis cs:0.556,1.66) {crosses};
  \draw[cHyd!55, line width=0.4pt] (axis cs:0.554,1.69) -- (axis cs:0.543,1.78);
  \node[ann, anchor=east, color=cPro] at (axis cs:0.630,3.30) {yeast ADH};
  \node[ann, anchor=west, color=cHyd] at (axis cs:0.424,1.09) {ecDHFR};

  % direction of falling commitment, inside the plot area
  \draw[-{Straight Barb[length=2.2pt,width=3.0pt]}, cGrey, line width=0.7pt]
    (axis cs:0.366,2.66) -- (axis cs:0.366,3.08);
  \node[annc, anchor=west, color=cGrey!80!black] at (axis cs:0.378,2.87)
    {masking\\increasing};
\end{axis}

% annotations placed outside the clipped axis
\node[ann, anchor=south west, rotate=54] at ($(A.south west)+(9mm,13.5mm)$)
  {semiclassical locus $X_{\mathrm H}=\gSC X_{\mathrm D}$};
\node[tag, fill=cBandG!20, anchor=north west] at ($(A.north west)+(2.0mm,-2.0mm)$)
  {$F>0$};
\node[tag, anchor=south east] at ($(A.south east)+(-2.0mm,2.0mm)$) {$F<0$};

% =====================================================================
% (B) the exponent moves along the ridge
% =====================================================================
\begin{axis}[geo, name=B, width=58mm, height=25mm,
  at={(A.north east)}, anchor=north west, xshift=21mm,
  xmode=log, xmin=1e-2, xmax=1.6e4, ymin=2.76, ymax=3.46,
  ylabel={exponent $X_{\mathrm H}/X_{\mathrm D}$},
  xticklabels={}, ytick={2.9,3.1,3.3}, clip=false,
  title={(\textit{\textbf{B}})\ \ along the ridge the exponent moves\ldots}]
  \addplot[cHyd, line width=1.2pt] table[x=s,y=expo] {data/fg_ridge.dat};
  \draw[cInk, line width=0.5pt, dashed] (axis cs:1e-2,3.34887) --
       (axis cs:1.6e4,3.34887);
  \node[ann, anchor=south east] at (axis cs:1.4e4,3.355) {$\gSC$};
  \node[ann, anchor=south east, color=cHyd] at (axis cs:1.3e4,2.80)
    {sweeps $(2.846,3.349)$};
\end{axis}

% =====================================================================
% (C) the offset does not
% =====================================================================
\begin{axis}[geo, name=C, width=58mm, height=25mm,
  at={(B.south west)}, anchor=north west, yshift=-11mm,
  xmode=log, xmin=1e-2, xmax=1.6e4, ymin=-0.115, ymax=0.035,
  xlabel={position along the gated ridge, $s=A/w^{2}$},
  ylabel={offset $F$}, ytick={-0.1,-0.042086,0},
  yticklabels={$-0.10$,$\Fz$,$0$}, clip=false,
  title={(\textit{\textbf{C}})\ \ \ldots and the offset does not}]
  \addplot[cPro, line width=1.2pt] table[x=s,y=off] {data/fg_ridge.dat};
  \draw[cInk, line width=0.5pt, dashed] (axis cs:1e-2,0) -- (axis cs:1.6e4,0);
  \node[ann, anchor=north east, color=cPro] at (axis cs:1.4e4,-0.050)
    {constant at $\Fz=-0.0421$};
\end{axis}
\end{tikzpicture}
\end{document}
